\documentclass[a4paper,14pt]{extarticle}

\usepackage[T2A]{fontenc}
\usepackage[utf8]{inputenc}
\usepackage[russian]{babel}
\usepackage[left=30mm,right=15mm,top=20mm,bottom=20mm]{geometry}
\usepackage{setspace}
\usepackage{indentfirst}
\usepackage{titlesec}
\usepackage{enumitem}
\usepackage{listings}
\usepackage{xcolor}
\usepackage{graphicx}
\usepackage{hyperref}
\usepackage{tabularx}
\usepackage{booktabs}
\usepackage{float}

\onehalfspacing
\setlength{\parindent}{1.25cm}

\definecolor{codegray}{rgb}{0.95,0.95,0.95}
\definecolor{codegreen}{rgb}{0,0.5,0}
\definecolor{codepurple}{rgb}{0.5,0,0.5}

\lstdefinestyle{json}{
  backgroundcolor=\color{codegray},
  basicstyle=\ttfamily\footnotesize,
  breaklines=true,
  captionpos=b,
  commentstyle=\color{codegreen},
  keywordstyle=\color{codepurple},
  stringstyle=\color{codegreen},
  frame=single,
  rulecolor=\color{gray},
  numbers=left,
  numberstyle=\tiny\color{gray},
  numbersep=5pt,
  tabsize=2,
}

\hypersetup{colorlinks=true, linkcolor=black, urlcolor=blue}

\begin{document}

% ===== ТИТУЛЬНЫЙ ЛИСТ =====
\begin{titlepage}
  \begin{center}
    \textbf{Министерство науки и высшего образования Российской Федерации}\\[0.3cm]
    \textbf{Федеральное государственное автономное образовательное учреждение\\
    высшего образования}\\[0.3cm]
    \textbf{«НАЦИОНАЛЬНЫЙ ИССЛЕДОВАТЕЛЬСКИЙ УНИВЕРСИТЕТ ИТМО»}\\[0.5cm]
    \hrule
    \vspace{0.3cm}
    Факультет программной инженерии и компьютерной техники\\[0.2cm]
    Образовательная программа «Разработка программного обеспечения»\\[2cm]

    \large\textbf{ОТЧЁТ}\\[0.3cm]
    \textbf{по домашнему заданию №3}\\[0.3cm]
    \textbf{по дисциплине «Бэкенд-разработка»}\\[0.5cm]
    \normalsize Тема: \textbf{Тестирование API средствами Postman}\\[3cm]

    \begin{flushright}
      \begin{tabular}{ll}
        Выполнил: & Иванов Виктор\\
        Группа:   & БР1.1\\
        Вариант:  & Система бронирования столиков\\
                  & в ресторанах\\[0.5cm]
        Проверил: & \\
      \end{tabular}
    \end{flushright}

    \vfill
    Санкт-Петербург, 2025
  \end{center}
\end{titlepage}

\tableofcontents
\newpage

% ===== 1. ЦЕЛЬ =====
\section{Цель работы}

Целью данного домашнего задания является реализация тестирования REST API приложения
бронирования столиков в ресторанах средствами Postman. В рамках работы необходимо:

\begin{itemize}[leftmargin=2cm]
  \item разработать комплексный сценарий тестирования по основному процессу приложения;
  \item реализовать автоматические тесты внутри Postman с помощью скриптов на JavaScript;
  \item обеспечить автоматическую передачу данных (токен, идентификаторы) между запросами;
  \item охватить как позитивные, так и негативные тестовые случаи.
\end{itemize}

% ===== 2. ЗАДАНИЕ =====
\section{Задание}

\begin{itemize}[leftmargin=2cm]
  \item Реализовать тестирование API средствами Postman.
  \item Написать тесты внутри Postman.
  \item Представить один комплексный сценарий по основному процессу приложения.
\end{itemize}

Основной сценарий для системы бронирования столиков:

\[
\textit{регистрация} \;\to\; \textit{вход} \;\to\; \textit{список ресторанов} \;\to\;
\textit{фильтрация} \;\to\; \textit{просмотр ресторана}
\]
\[
\;\to\; \textit{проверка доступности} \;\to\; \textit{бронирование} \;\to\;
\textit{отзыв} \;\to\; \textit{отмена брони}
\]

% ===== 3. РЕАЛИЗОВАННЫЙ API =====
\section{Реализованный API}

В качестве тестируемого сервиса выступает REST API, разработанный в рамках лабораторной
работы №1. Стек технологий:

\begin{itemize}[leftmargin=2cm]
  \item \textbf{Runtime:} Node.js
  \item \textbf{Framework:} Express.js
  \item \textbf{ORM:} TypeORM
  \item \textbf{База данных:} SQLite
  \item \textbf{Аутентификация:} JWT (Bearer токен)
  \item \textbf{Документация:} Swagger UI (\texttt{/api-docs})
\end{itemize}

API запущен локально по адресу \texttt{http://localhost:3001}.
Базовый префикс всех эндпоинтов: \texttt{/api/v1}.

\subsection{Реализованные эндпоинты}

\begin{table}[H]
\centering
\caption{Эндпоинты REST API}
\begin{tabularx}{\textwidth}{llX}
\toprule
\textbf{Метод} & \textbf{Путь} & \textbf{Описание} \\
\midrule
POST   & /auth/register                       & Регистрация пользователя \\
POST   & /auth/login                          & Вход, получение JWT \\
GET    & /users/me                            & Профиль пользователя \\
PATCH  & /users/me                            & Обновление профиля \\
PUT    & /users/me/password                   & Смена пароля \\
GET    & /restaurants                         & Список ресторанов \\
GET    & /restaurants/\{id\}                  & Детали ресторана \\
GET    & /restaurants/\{id\}/tables           & Столики ресторана \\
GET    & /restaurants/\{id\}/availability     & Доступность столиков \\
GET    & /restaurants/\{id\}/reviews          & Отзывы о ресторане \\
POST   & /restaurants/\{id\}/reviews          & Добавить отзыв \\
GET    & /bookings                            & Список броней пользователя \\
POST   & /bookings                            & Создать бронирование \\
GET    & /bookings/\{id\}                     & Детали бронирования \\
DELETE & /bookings/\{id\}                     & Отменить бронирование \\
\bottomrule
\end{tabularx}
\end{table}

% ===== 4. ТЕСТОВЫЙ СЦЕНАРИЙ =====
\section{Тестовый сценарий}

\subsection{Структура коллекции}

Postman-коллекция \textbf{«Restaurant Booking API»} содержит 18 запросов,
сгруппированных в один сквозной сценарий. Для автоматической передачи данных между
шагами используются переменные коллекции:

\begin{table}[H]
\centering
\caption{Переменные коллекции}
\begin{tabularx}{\textwidth}{lX}
\toprule
\textbf{Переменная} & \textbf{Назначение} \\
\midrule
\texttt{base\_url}      & Базовый адрес API (\texttt{http://localhost:3001/api/v1}) \\
\texttt{access\_token}  & JWT токен, сохраняется после регистрации/логина \\
\texttt{restaurant\_id} & ID первого ресторана из списка \\
\texttt{table\_id}      & ID первого столика выбранного ресторана \\
\texttt{booking\_id}    & ID созданного бронирования \\
\texttt{test\_email}    & Уникальный email (генерируется в pre-request скрипте) \\
\bottomrule
\end{tabularx}
\end{table}

\subsection{Шаги сценария}

\begin{enumerate}[leftmargin=2cm]
  \item \textbf{Регистрация нового пользователя} — \texttt{POST /auth/register}.\\
    Pre-request скрипт генерирует уникальный email на основе timestamp.
    Токен сохраняется в \texttt{access\_token}.

  \item \textbf{Вход в систему} — \texttt{POST /auth/login}.\\
    Проверка что логин возвращает валидный JWT.

  \item \textbf{Профиль пользователя} — \texttt{GET /users/me}.\\
    Проверка авторизации через Bearer токен.

  \item \textbf{Список ресторанов} — \texttt{GET /restaurants}.\\
    Сохранение \texttt{restaurant\_id} из первого элемента ответа.

  \item \textbf{Фильтрация ресторанов} — \texttt{GET /restaurants?cuisine\_type=русская}.\\
    Проверка что все возвращённые рестораны соответствуют фильтру.

  \item \textbf{Детали ресторана} — \texttt{GET /restaurants/\{restaurant\_id\}}.\\
    Проверка полноты данных и наличия счётчиков отзывов.

  \item \textbf{Список столиков} — \texttt{GET /restaurants/\{id\}/tables}.\\
    Сохранение \texttt{table\_id} для создания бронирования.

  \item \textbf{Проверка доступности} — \texttt{GET /restaurants/\{id\}/availability}.\\
    Запрос с параметрами \texttt{date}, \texttt{time}, \texttt{guests}.

  \item \textbf{Создание бронирования} — \texttt{POST /bookings}.\\
    Сохранение \texttt{booking\_id}. Проверка статуса \texttt{confirmed}.

  \item \textbf{Детали бронирования} — \texttt{GET /bookings/\{booking\_id\}}.\\
    Проверка вложенных данных ресторана и столика.

  \item \textbf{Список бронирований} — \texttt{GET /bookings?status=confirmed}.\\
    Проверка фильтрации по статусу.

  \item \textbf{Добавление отзыва} — \texttt{POST /restaurants/\{id\}/reviews}.\\
    Оценка 5, комментарий. Проверка данных автора в ответе.

  \item \textbf{Просмотр отзывов} — \texttt{GET /restaurants/\{id\}/reviews}.\\
    Проверка что отзыв появился в списке.

  \item \textbf{Отмена бронирования} — \texttt{DELETE /bookings/\{booking\_id\}}.\\
    Ожидаемый статус \texttt{204 No Content}.

  \item \textbf{Повторная отмена (негативный тест)} — \texttt{DELETE /bookings/\{booking\_id\}}.\\
    Ожидается \texttt{409 Conflict} с кодом \texttt{cannot\_cancel\_booking}.

  \item \textbf{Запрос без токена (негативный тест)} — \texttt{GET /users/me}.\\
    Ожидается \texttt{401 Unauthorized}.

  \item \textbf{Неверный пароль (негативный тест)} — \texttt{POST /auth/login}.\\
    Ожидается \texttt{401} с кодом \texttt{invalid\_credentials}.

  \item \textbf{Несуществующий ресторан (негативный тест)} — \texttt{GET /restaurants/00000000-...}.\\
    Ожидается \texttt{404 Not Found}.
\end{enumerate}

\subsection{Пример теста}

Ниже приведён пример тестового скрипта для запроса «Создание бронирования»:

\begin{lstlisting}[style=json, caption={Тест для POST /bookings}]
pm.test('Статус 201 - бронирование создано', function () {
    pm.response.to.have.status(201);
});

pm.test('Статус бронирования = confirmed', function () {
    pm.expect(pm.response.json().status).to.equal('confirmed');
});

pm.test('Бронирование включает данные ресторана и столика', function () {
    var json = pm.response.json();
    pm.expect(json).to.have.property('restaurant');
    pm.expect(json).to.have.property('table');
});

// Сохраняем booking_id для следующих шагов
var json = pm.response.json();
if (json.id) {
    pm.collectionVariables.set('booking_id', json.id);
}
\end{lstlisting}

% ===== 5. РЕЗУЛЬТАТЫ =====
\section{Результаты тестирования}

Коллекция запущена через Collection Runner Postman (Run collection).

\begin{table}[H]
\centering
\caption{Итоги прогона коллекции}
\begin{tabular}{lr}
\toprule
\textbf{Показатель} & \textbf{Значение} \\
\midrule
Всего запросов       & 18 \\
Всего тестов         & 55 \\
Passed               & \textbf{55} \\
Failed               & 0 \\
Errors               & 0 \\
Skipped              & 0 \\
Среднее время ответа & 17 мс \\
Общая длительность   & 827 мс \\
\bottomrule
\end{tabular}
\end{table}

Все 55 проверок прошли успешно. Как позитивные сценарии (регистрация, бронирование,
отзыв), так и негативные (401, 404, 409) возвращают ожидаемые коды и тела ответов.

% ===== 6. ВЫВОД =====
\section{Вывод}

В ходе выполнения домашнего задания был реализован комплексный сценарий тестирования
REST API системы бронирования столиков в ресторанах средствами Postman. Сценарий
покрывает полный жизненный цикл основного бизнес-процесса: от регистрации пользователя
до создания и отмены бронирования.

Ключевые результаты:
\begin{itemize}[leftmargin=2cm]
  \item разработана Postman-коллекция из 18 запросов с 55 автоматическими проверками;
  \item реализована автоматическая передача токенов и идентификаторов между шагами
        через переменные коллекции;
  \item все тесты успешно пройдены (\textbf{55/55 Passed, 0 Failed});
  \item покрыты негативные сценарии: попытка доступа без токена (401),
        повторная отмена брони (409), запрос несуществующего ресурса (404).
\end{itemize}

\end{document}
